\section{Discussion}

Research on international scholarly mobility has documented a widening range of origins and destinations \cite{ref008,yenilmez2026}, while studies of scientific cities have identified persistent differences in their positions within researcher mobility networks \cite{ref031,ref032}. These findings raise a question about how geographic inequalities change as researcher mobility expands. Our results reveal two distinct patterns. International mobility became more evenly distributed across countries, yet substantial increases in worldwide shares were concentrated in a relatively small number of countries. Within countries, the growing number of cities with little recorded mobility contributed to widening relative differences across an expanding range of cities, even as flows became more evenly distributed among cities recorded every year in many countries. International diversification and widening relative differences across cities therefore describe different aspects of geographic change: the former occurred alongside unequal gains in national shares, while the latter was strongly influenced by the growing presence of cities with small flows.

The growing diversity of international researcher mobility has been documented through changes in the countries researchers move from and the destinations they choose \cite{ref008,ref039}. Studies of global mobility networks have also identified unequal positions among countries associated with differences in scientific capacity \cite{ref009}. More recently, Yenilmez \cite{yenilmez2026} found that international academic migration has become less concentrated around its largest hubs, while the United States remains central and several Asian countries have gained prominence. Our findings support this broader picture of international diversification but reveal substantial differences in how countries' shares of worldwide movement have changed. The United States retained the largest share despite its relative decline, while China and India gained considerably. Although many countries increased their shares of worldwide researcher mobility, most of the combined gains were concentrated among relatively few countries. International diversification therefore involved a more even global distribution of mobility through highly unequal changes in national shares. Established research countries remained prominent, while a relatively small group of countries recorded much of the increase in worldwide shares.

Research on the geography of science has documented the declining concentration of scientific publications in major cities \cite{ref015} and the growth of scientific collaboration between cities, particularly within countries \cite{ref016}. These findings show how the geography of scientific activity has changed, but do not directly establish how researcher mobility has evolved across cities. Studies of researcher mobility have shown that major scientific cities attract and retain highly productive researchers \cite{ref031}, while substantial differences persist between central and peripheral cities in mobility networks \cite{ref032}. Our findings show that domestic researcher mobility was recorded in an increasing number of cities, although many cities first appearing in later years accounted for relatively little movement. Taken together, cities with mobility recorded every year continued to account for most domestic flows. In many countries, however, flows became more evenly distributed among these cities, even as relative differences across the wider range of cities increased. Further analysis showed that cities first appearing in the records after 1998 made a substantial contribution to the rise in relative inequality measured across cities each year. This changes how that rise should be understood: it can reflect the appearance of many cities with small flows, even as differences among cities recorded each year narrow.

Differences in researcher mobility across scientific disciplines have been documented in studies of mobility patterns and career stages \cite{ref033}, international movements among postdoctoral researchers \cite{ref035}, and discipline-specific talent gains and losses \cite{ref034}. These studies establish important differences in researchers' movements and career experiences, raising a related question about how mobility is distributed geographically within each field. Our findings show that physical, life, health, and social sciences shared a broad pattern: international flows became more evenly distributed across countries, while relative inequality across the growing number of cities recorded each year increased in many countries. The increase was particularly pronounced in social sciences. In a separate comparison of the same 34 countries across all four disciplines, however, the median domestic trend shifted toward a more even distribution when only cities with mobility recorded each year were compared. This occurred in every discipline, including social sciences. In this matched-country comparison, the increase observed across the expanding range of cities was not accompanied by a similar increase in inequality among the same cities. The distinction found in the overall domestic analysis also applies across scientific fields, even though the magnitude of the observed changes differs between disciplines.

The interpretation of these findings depends on the movements captured in publication records and on changes in the availability of those records. Comparisons using similar numbers of recorded movements across years substantially reduced the number of countries meeting the criteria for concentration in domestic mobility. This sensitivity matters because the number of recorded movements increased considerably over the study period. Bibliometric studies have also shown that changes in publication affiliations can reflect different forms of academic mobility, including migration, temporary appointments and multiple institutional affiliations \cite{ref022,ref004}. A city's first appearance in publication records therefore does not establish when researchers first began moving to or from that city. Differences in publication coverage, institutional addresses and geographic identification can further affect long-term comparisons. Our domestic analysis is restricted to countries with sufficiently consistent records, and the disciplinary comparisons cover a smaller group of countries. Finally, the combined volume of incoming and outgoing movement measures the extent of involvement in researcher mobility, while assessments of net talent gains, losses and retention require additional information \cite{ref041,ref034}. Changes in publication coverage and recorded event volume are therefore particularly important when interpreting long-term trends in the geographic distribution of researcher mobility.

The findings raise further questions about how the geography of researcher mobility develops over time. One important direction is to follow cities that first appear in publication records with relatively little movement. Studies of subnational researcher migration have identified substantial regional differences in attracting scientific talent \cite{ref029}. Examining subsequent movements to and from these cities could show whether initially small flows develop into sustained exchanges, whether cities gain or lose researchers over time, and whether their shares of domestic mobility increase. At the international level, changes in countries' worldwide mobility shares invite further comparisons with scientific employment, research output and research influence. Studies linking mobility to knowledge transfer and the development of national science systems provide a basis for examining these relationships \cite{ref006,ref040}. Further research could also examine how periods of growth, slower expansion and decline across scientific disciplines relate to changes in the geography of movement. Understanding long-term geographic inequality will require following both where researcher mobility is recorded and how the shares of movement change among places already represented in those records.
